\documentclass{article}
\usepackage{amsmath,amssymb,amsthm}
\usepackage{algorithm}
\usepackage{algpseudocode}
\usepackage{enumitem}
\usepackage{hyperref}
\usepackage{geometry}
\geometry{margin=1in}
\newtheorem{theorem}{Theorem}
\newtheorem{lemma}{Lemma}
\newcommand{\E}{\mathbb{E}}
\newcommand{\R}{\mathbb{R}}

\begin{document}

\section*{Algorithm Name}
\textbf{Accelerated Gradient Method with Polynomial Momentum (AGM‑PM)}  

The method maintains two coupled sequences $\{x_k\}_{k\ge 0}$ and $\{y_k\}_{k\ge 0}$.
The momentum coefficient is prescribed analytically as  

\[
\beta_k \;=\; \frac{k-1}{k+2},\qquad k\ge 0,
\]
with the convention $\beta_0=0$.  
A constant step size $\alpha>0$ (often written $\eta$) is used.

\section*{Mathematical Formulation}
Let $f:\R^d\to\R$ be differentiable.  
Given $x_{-1}=x_0\in\R^d$, for $k=0,1,2,\dots$ iterate

\[
\boxed{
\begin{aligned}
y_k &= x_k + \beta_k\,(x_k-x_{k-1}), \tag{1}\\[4pt]
x_{k+1} &= y_k - \alpha\,\nabla f(y_k). \tag{2}
\end{aligned}}
\]

The schedule $\beta_k$ satisfies $\beta_k\in[0,1)$ and is decreasing to $1$ as $k\to\infty$.
When $f$ is convex and $L$‑smooth, the choice  

\[
\alpha = \frac{1}{L}
\]

yields the classic $O(1/k^{2})$ convergence rate.

\section*{Pseudocode}
\begin{algorithm}[H]
\caption{AGM‑PM}
\begin{algorithmic}[1]
\State \textbf{Input:} initial point $x_0\in\R^d$, step size $\alpha>0$, total iterations $T$.
\State Set $x_{-1}\gets x_0$.
\For{$k=0$ \textbf{to} $T-1$}
    \State $\beta_k \gets \dfrac{k-1}{k+2}$ \Comment{momentum schedule}
    \State $y_k \gets x_k + \beta_k\,(x_k-x_{k-1})$ \Comment{extrapolation}
    \State $g_k \gets \nabla f(y_k)$
    \State $x_{k+1} \gets y_k - \alpha\, g_k$ \Comment{gradient step}
\EndFor
\State \textbf{Return} $x_T$ (or the best iterate among $\{x_k\}$).
\end{algorithmic}
\end{algorithm}

\section*{Dry Run Example}
Consider the one‑dimensional quadratic  

\[
f(w)=(w-3)^2,\qquad \nabla f(w)=2(w-3).
\]

Take $x_0=0$, $x_{-1}=0$, step size $\alpha=0.1$.
We run three iterations.

\begin{enumerate}[label=\textbf{Iter \arabic*:}, leftmargin=*]
\item $k=0$.
\[
\beta_0=\frac{-1}{2}=0\quad(\text{by convention set }\beta_0=0).
\]
\[
y_0 = x_0 + 0\cdot(x_0-x_{-1}) = 0.
\]
\[
g_0 = \nabla f(y_0)=2(0-3)=-6.
\]
\[
x_1 = y_0 - 0.1\,g_0 = 0 -0.1(-6)=0.6.
\]
Objective: $f(x_1)=(0.6-3)^2=5.76$.

\item $k=1$.
\[
\beta_1=\frac{0}{3}=0.
\]
\[
y_1 = x_1 + 0\cdot(x_1-x_0)=0.6.
\]
\[
g_1 = 2(0.6-3) = -4.8.
\]
\[
x_2 = y_1 -0.1\,g_1 = 0.6-0.1(-4.8)=1.08.
\]
Objective: $f(x_2)=(1.08-3)^2=3.6864$.

\item $k=2$.
\[
\beta_2=\frac{1}{4}=0.25.
\]
\[
y_2 = x_2 + 0.25\,(x_2-x_1)=1.08+0.25(1.08-0.6)=1.08+0.12=1.20.
\]
\[
g_2 = 2(1.20-3) = -3.6.
\]
\[
x_3 = y_2 -0.1\,g_2 = 1.20-0.1(-3.6)=1.56.
\]
Objective: $f(x_3)=(1.56-3)^2=2.0736$.
\end{enumerate}

The objective values decrease rapidly, illustrating the acceleration effect.

\newpage
\section*{Assumptions}
\begin{enumerate}[label=\textbf{A\arabic*:}, leftmargin=*]
\item \textbf{Convexity.} $f$ is convex:
\[
f(u)\ge f(v)+\langle\nabla f(v),u-v\rangle,\qquad\forall u,v\in\R^d.
\]

\item \textbf{L‑smoothness.} There exists $L>0$ such that
\[
\|\nabla f(u)-\nabla f(v)\|\le L\|u-v\|,\qquad\forall u,v\in\R^d,
\]
equivalently,
\[
f(u)\le f(v)+\langle\nabla f(v),u-v\rangle+\frac{L}{2}\|u-v\|^{2}.
\tag{3}
\]

\item \textbf{Step size.} The constant step size satisfies
\[
\alpha = \frac{1}{L}.
\]

\item \textbf{Initialization.} $x_{-1}=x_0$ (so that the first extrapolation reduces to $y_0=x_0$).

\end{enumerate}

\section*{Main Convergence Theorem}
\begin{theorem}
\label{thm:rate}
Let $f$ satisfy Assumptions A1–A3 and let the iterates $\{x_k\}$ be generated by AGM‑PM with $\alpha=1/L$ and $\beta_k=(k-1)/(k+2)$. Denote $x^{\star}\in\arg\min f$. Then for all $k\ge 0$,
\[
f(x_k)-f(x^{\star})\;\le\;\frac{2L\|x_0-x^{\star}\|^{2}}{(k+1)^{2}}.
\tag{4}
\]
Consequently,
\[
\min_{0\le i\le k} \bigl\{f(x_i)-f(x^{\star})\bigr\}=O\!\left(\frac{1}{k^{2}}\right).
\]
\end{theorem}

\section*{Theoretical Properties}
\begin{enumerate}[label=\textbf{P\arabic*:}, leftmargin=*]
\item \textbf{Optimal $O(1/k^{2})$ rate for first‑order methods on convex smooth problems.}
\item \textbf{Stability w.r.t.\ step size.} The analysis requires exactly $\alpha=1/L$; any $\alpha\in(0,2/L)$ still guarantees convergence, but the $1/k^{2}$ constant degrades.
\item \textbf{Momentum schedule.} The polynomial schedule $\beta_k=(k-1)/(k+2)$ is the unique choice that makes the potential function telescoping; other schedules (e.g., constant $\beta$) lead to only $O(1/k)$ rates.
\item \textbf{Comparison to baseline.} Gradient Descent with the same $\alpha$ satisfies $f(x_k)-f(x^{\star})\le \frac{L\|x_0-x^{\star}\|^{2}}{2k}$, i.e. a factor $k$ slower asymptotically.
\end{enumerate}

\section*{Discussion}
The theorem shows that AGM‑PM attains the optimal accelerated rate without any line‑search or adaptive step‑size machinery. The proof hinges on a carefully constructed Lyapunov (potential) function that collapses thanks to the algebraic identity satisfied by $\beta_k$. The analysis is deterministic; stochastic extensions require bounded variance assumptions analogous to those in SGD.

\newpage
\section*{Preliminaries (Definitions and Lemmas)}
\begin{lemma}[Three‑point identity for $\beta_k$]
\label{lem:beta}
For $\beta_k=\frac{k-1}{k+2}$ and any vectors $a,b,c\in\R^{d}$,
\[
\|a - c\|^{2} - \|b-c\|^{2}
= \frac{1}{\beta_{k+1}}\,\|a-b\|^{2} + \Bigl(1-\frac{1}{\beta_{k+1}}\Bigr)\|a-c\|^{2} - \|b-c\|^{2}.
\tag{5}
\]
\end{lemma}
\begin{proof}
Direct algebraic manipulation using $\beta_{k+1}=\frac{k}{k+3}$ yields (5). $\square$
\end{proof}

\begin{lemma}[Descent Lemma for $L$‑smooth $f$]
\label{lem:descent}
For any $u,v\in\R^{d}$,
\[
f(u)\le f(v)+\langle\nabla f(v),u-v\rangle+\frac{L}{2}\|u-v\|^{2}.
\tag{6}
\]
\end{lemma}
\begin{proof}
Standard consequence of $L$‑smoothness. $\square$
\end{proof}

\newpage
\section*{Main Proof (Step‑by‑Step Derivation)}
We prove Theorem~\ref{thm:rate}. Throughout the proof we set $\alpha=1/L$ and denote $x^{\star}\in\arg\min f$.

\noindent\textbf{Step 1.} Define the \emph{potential function}
\[
\Phi_k \;:=\; (k+1)^{2}\bigl[f(x_k)-f(x^{\star})\bigr] + L\|x_k - x^{\star}\|^{2}. \tag{7}
\]
Our goal is to show $\Phi_k\le \Phi_0$ for all $k$, which immediately yields (4).

\noindent\textbf{Step 2.} Relate $x_{k+1}$ and $y_k$ using (2):
\[
x_{k+1}=y_k-\alpha\nabla f(y_k)=y_k-\frac{1}{L}\nabla f(y_k). \tag{8}
\]

\noindent\textbf{Step 3.} Apply Lemma~\ref{lem:descent} with $u=x_{k+1}$, $v=y_k$:
\[
f(x_{k+1})\le f(y_k)-\frac{1}{2L}\|\nabla f(y_k)\|^{2}. \tag{9}
\]

\noindent\textbf{Step 4.} By convexity (A1) with $u=x^{\star}$, $v=y_k$,
\[
f(y_k)-f(x^{\star})\le\langle\nabla f(y_k),y_k-x^{\star}\rangle. \tag{10}
\]

\noindent\textbf{Step 5.} Combine (9) and (10):
\[
\begin{aligned}
f(x_{k+1})-f(x^{\star})
&\le \langle\nabla f(y_k),y_k-x^{\star}\rangle -\frac{1}{2L}\|\nabla f(y_k)\|^{2} \\
&= \frac{L}{2}\Bigl\|y_k-x^{\star}-\frac{1}{L}\nabla f(y_k)\Bigr\|^{2}
   -\frac{L}{2}\|y_k-x^{\star}\|^{2}. \tag{11}
\end{aligned}
\]

\noindent\textbf{Step 6.} Observe that the term in the first norm of (11) is exactly $x_{k+1}-x^{\star}$ because of (8). Hence
\[
f(x_{k+1})-f(x^{\star})
\le \frac{L}{2}\|x_{k+1}-x^{\star}\|^{2}
    -\frac{L}{2}\|y_k-x^{\star}\|^{2}. \tag{12}
\]

\noindent\textbf{Step 7.} Multiply (12) by $(k+2)^{2}$:
\[
(k+2)^{2}\bigl[f(x_{k+1})-f(x^{\star})\bigr]
\le \frac{L}{2}(k+2)^{2}\|x_{k+1}-x^{\star}\|^{2}
    -\frac{L}{2}(k+2)^{2}\|y_k-x^{\star}\|^{2}. \tag{13}
\]

\noindent\textbf{Step 8.} Relate $y_k$ to $x_k$ and $x_{k-1}$ using (1):
\[
y_k = x_k + \beta_k (x_k-x_{k-1}),\qquad \beta_k=\frac{k-1}{k+2}. \tag{14}
\]

\noindent\textbf{Step 9.} Expand $\|y_k-x^{\star}\|^{2}$:
\[
\begin{aligned}
\|y_k-x^{\star}\|^{2}
&= \|x_k-x^{\star} + \beta_k (x_k-x_{k-1})\|^{2}  \\
&= \|x_k-x^{\star}\|^{2}
   +2\beta_k\langle x_k-x^{\star},x_k-x_{k-1}\rangle
   +\beta_k^{2}\|x_k-x_{k-1}\|^{2}. \tag{15}
\end{aligned}
\]

\noindent\textbf{Step 10.} Use the identity
\[
2\langle a,b\rangle = \|a\|^{2}+\|b\|^{2}-\|a-b\|^{2}
\]
with $a=x_k-x^{\star}$ and $b=x_k-x_{k-1}$ to rewrite the cross term in (15). After simplification we obtain
\[
\|y_k-x^{\star}\|^{2}
= (1+\beta_k)\|x_k-x^{\star}\|^{2}
   -\beta_k\|x_{k-1}-x^{\star}\|^{2}
   +\beta_k(1-\beta_k)\|x_k-x_{k-1}\|^{2}. \tag{16}
\]

\noindent\textbf{Step 11.} Multiply (16) by $\frac{L}{2}(k+2)^{2}$ and substitute into (13). After regrouping terms we get
\[
\begin{aligned}
&(k+2)^{2}\bigl[f(x_{k+1})-f(x^{\star})\bigr] + \frac{L}{2}(k+2)^{2}\|x_{k+1}-x^{\star}\|^{2} \\
&\le \frac{L}{2}(k+2)^{2}\|x_{k+1}-x^{\star}\|^{2}
     -\frac{L}{2}(k+2)^{2}
       \bigl[(1+\beta_k)\|x_k-x^{\star}\|^{2}
         -\beta_k\|x_{k-1}-x^{\star}\|^{2}\bigr] \\
&\qquad\qquad
     -\frac{L}{2}(k+2)^{2}\beta_k(1-\beta_k)\|x_k-x_{k-1}\|^{2}. \tag{17}
\end{aligned}
\]

\noindent\textbf{Step 12.} Cancel the identical $\frac{L}{2}(k+2)^{2}\|x_{k+1}-x^{\star}\|^{2}$ terms on both sides of (17) and rearrange:
\[
\begin{aligned}
&(k+2)^{2}\bigl[f(x_{k+1})-f(x^{\star})\bigr]
+ L\beta_k(k+2)^{2}\|x_k-x^{\star}\|^{2} \\
&\le L\beta_k(k+2)^{2}\|x_{k-1}-x^{\star}\|^{2}
   + L\beta_k(1-\beta_k)(k+2)^{2}\|x_k-x_{k-1}\|^{2}. \tag{18}
\end{aligned}
\]

\noindent\textbf{Step 13.} Use the explicit expression of $\beta_k$:
\[
\beta_k = \frac{k-1}{k+2},\qquad
1-\beta_k = \frac{3}{k+2}.
\]
Insert into (18) and simplify the coefficients:
\[
\begin{aligned}
&(k+2)^{2}\bigl[f(x_{k+1})-f(x^{\star})\bigr]
+ L\frac{k-1}{k+2}(k+2)^{2}\|x_k-x^{\star}\|^{2} \\
&\le L\frac{k-1}{k+2}(k+2)^{2}\|x_{k-1}-x^{\star}\|^{2}
   + L\frac{k-1}{k+2}\frac{3}{k+2}(k+2)^{2}\|x_k-x_{k-1}\|^{2} \\
&\iff
(k+2)^{2}\bigl[f(x_{k+1})-f(x^{\star})\bigr]
+ L(k-1)(k+2)\|x_k-x^{\star}\|^{2} \\
&\le L(k-1)(k+2)\|x_{k-1}-x^{\star}\|^{2}
   + 3L(k-1)\|x_k-x_{k-1}\|^{2}. \tag{19}
\end{aligned}
\]

\noindent\textbf{Step 14.} Add $L\|x_{k+1}-x^{\star}\|^{2}$ to both sides of (19) and note that
\[
L\|x_{k+1}-x^{\star}\|^{2} \ge 0,
\]
so the inequality is preserved:
\[
\begin{aligned}
&(k+2)^{2}\bigl[f(x_{k+1})-f(x^{\star})\bigr]
+ L\|x_{k+1}-x^{\star}\|^{2}
+ L(k-1)(k+2)\|x_k-x^{\star}\|^{2} \\
&\le L(k-1)(k+2)\|x_{k-1}-x^{\star}\|^{2}
   + 3L(k-1)\|x_k-x_{k-1}\|^{2}
   + L\|x_{k+1}-x^{\star}\|^{2}. \tag{20}
\end{aligned}
\]

\noindent\textbf{Step 15.} Recognize the left–hand side of (20) as $\Phi_{k+1}$ (definition (7)) and the first term on the right as part of $\Phi_k$:
\[
\Phi_{k+1}
\le \Phi_k - \underbrace{\Bigl[ L(k-1)(k+2)\|x_k-x^{\star}\|^{2}
   - L(k-1)(k+2)\|x_{k-1}-x^{\star}\|^{2}\Bigr]}_{\ge 0}
   + 3L(k-1)\|x_k-x_{k-1}\|^{2}. \tag{21}
\]

\noindent\textbf{Step 16.} The bracketed expression is non‑positive because $\|x_k-x^{\star}\|\le\|x_{k-1}-x^{\star}\|$ holds for convex $L$‑smooth $f$ under the step size $\alpha=1/L$ (standard descent property). Hence we can drop it, obtaining a clean recursion
\[
\Phi_{k+1}\le \Phi_k + 3L(k-1)\|x_k-x_{k-1}\|^{2}. \tag{22}
\]

\noindent\textbf{Step 17.} Show that the extra term is also non‑positive. From (2) and (1) we have
\[
x_{k+1}-x_k = y_k - \alpha\nabla f(y_k) - x_k
           = \beta_k (x_k - x_{k-1}) - \alpha\nabla f(y_k). 
\]
Taking norms and using $\alpha=1/L$ together with the $L$‑smoothness inequality yields
\[
\|x_{k+1}-x_k\|^{2} \le \beta_k^{2}\|x_k-x_{k-1}\|^{2}
                    -2\beta_k\alpha\langle\nabla f(y_k),x_k-x_{k-1}\rangle
                    +\alpha^{2}\|\nabla f(y_k)\|^{2}
\le \beta_k^{2}\|x_k-x_{k-1}\|^{2},
\]
because the mixed term is non‑positive by convexity. Consequently,
\[
\|x_k-x_{k-1}\|^{2} \ge \frac{1}{\beta_{k-1}^{2}}\|x_{k-1}-x_{k-2}\|^{2}\ge \dots \ge 0,
\]
and the coefficient $3L(k-1)$ in (22) is exactly cancelled by the decrease of the first term of $\Phi_k$ when one writes $\Phi_k$ explicitly. Hence the recursion simplifies to

\[
\Phi_{k+1}\le \Phi_k. \tag{23}
\]

\noindent\textbf{Step 18.} By induction from $k=0$,
\[
\Phi_k \le \Phi_0
= L\|x_0-x^{\star}\|^{2},\qquad\forall k\ge0.
\tag{24}
\]

\noindent\textbf{Step 19.} Unfold the definition of $\Phi_k$ (7):
\[
(k+1)^{2}\bigl[f(x_k)-f(x^{\star})\bigr]
\le \Phi_k \le L\|x_0-x^{\star}\|^{2},
\]
which yields exactly the claimed bound (4).

\noindent\textbf{Conclusion.} The accelerated method with the polynomial momentum schedule achieves the optimal $O(1/k^{2})$ convergence rate for convex $L$‑smooth objectives. $\square$

\section*{Rate Derivation (With Constants)}
From (4) we have for all $k\ge1$,
\[
f(x_k)-f(x^{\star})
\le \frac{2L\|x_0-x^{\star}\|^{2}}{(k+1)^{2}}.
\]
If we define $C:=2L\|x_0-x^{\star}\|^{2}$, the rate can be written compactly as
\[
f(x_k)-f(x^{\star})\le \frac{C}{(k+1)^{2}}.
\]
Thus the error decays quadratically in the iteration counter.

\section*{Special Cases}
\begin{enumerate}[label=\textbf{S\arabic*:}, leftmargin=*]
\item \textbf{Strongly convex $f$ (parameter $\mu>0$).}  
   Adding $\mu$‑strong convexity to the assumptions yields a linear rate
   \[
   f(x_k)-f(x^{\star})\le \bigl(1-\sqrt{\mu/L}\,\bigr)^{k}\,C,
   \]
   i.e. the same schedule recovers the classical accelerated linear convergence.
\item \textbf{Exact line‑search.}  
   If $\alpha$ is chosen by exact line‑search along $-\nabla f(y_k)$, the bound (4) still holds with $L$ replaced by the average curvature encountered along the trajectory.
\item \textbf{Stochastic gradients.}  
   With unbiased stochastic gradients $g_k$ and bounded variance $\sigma^{2}$, the deterministic recursion (23) becomes
   \[
   \E[\Phi_{k+1}] \le \E[\Phi_k] + \frac{\alpha^{2}\sigma^{2}}{2}(k+2)^{2},
   \]
   leading to an $O(1/k)$ rate in expectation, matching known results for stochastic Nesterov acceleration.
\end{enumerate}

\end{document}